% --- Archon physics formalization source begin ---
% archon:physics
% archon:covers PhyXMiniProblems/problem_phyx_mini_0963.lean
% archon:source-report reports/phyx_mini/problem_phyx_mini_0963.source.json

\chapter{Physics problem phyx\_mini\_0963}
\label{ch:PhyXMiniProblems_problem_phyx_mini_0963}

\paragraph{Problem source.}
\#\# Physical scenario

A pair of point charges, \textbackslash{}( q = +8.00\textbackslash{},\textbackslash{}mu\textbackslash{}text\{C\} \textbackslash{}) and \textbackslash{}( q' = -5.00\textbackslash{},\textbackslash{}mu\textbackslash{}text\{C\} \textbackslash{}), are moving with speeds \textbackslash{}( v = 9.00 \textbackslash{}times 10\textasciicircum{}4\textbackslash{},\textbackslash{}text\{m/s\} \textbackslash{}) and \textbackslash{}( v' = 6.50 \textbackslash{}times 10\textasciicircum{}4\textbackslash{},\textbackslash{}text\{m/s\} \textbackslash{}).

\#\# Question

When the charges are at the positions shown, find the magnetic force that \textbackslash{}( q' \textbackslash{}) exerts on \textbackslash{}( q \textbackslash{}).

\#\# Answer choices

- A: A: (7.49 \textbackslash{}times 10\textasciicircum{}\{-8\}N)\textbackslash{}hat\{\textbackslash{}mathbf\{i\}\}

- B: B: (7.49 \textbackslash{}times 10\textasciicircum{}\{-8\}N)\textbackslash{}hat\{\textbackslash{}mathbf\{j\}\}

- C: C: (7.49 \textbackslash{}times 10\textasciicircum{}\{-9\}N)\textbackslash{}hat\{\textbackslash{}mathbf\{j\}\}

- D: D: (1.50 \textbackslash{}times 10\textasciicircum{}\{-8\}N)\textbackslash{}hat\{\textbackslash{}mathbf\{j\}\}

\#\# Figure caption (auxiliary; use the image as primary evidence)

The image shows a coordinate system with the origin labeled \textbackslash{}(O\textbackslash{}). There are two charged particles depicted:

1. **Particle \textbackslash{}(q\textbackslash{}):** 
   - Positioned at \textbackslash{}(0.300 \textbackslash{}, \textbackslash{}text\{m\}\textbackslash{}) on the \textbackslash{}(y\textbackslash{})-axis.
   - Represented by a red circle with a positive sign inside it.
   - An arrow labeled \textbackslash{}(v\textbackslash{}) points horizontally to the right, indicating velocity.

2. **Particle \textbackslash{}(q'\textbackslash{}):**
   - Positioned at \textbackslash{}(0.400 \textbackslash{}, \textbackslash{}text\{m\}\textbackslash{}) on the \textbackslash{}(x\textbackslash{})-axis.
   - Represented by a blue circle with a negative sign inside it.
   - An arrow labeled \textbackslash{}(v'\textbackslash{}) points vertically upwards, indicating velocity.

The \textbackslash{}(x\textbackslash{})-axis is marked with \textbackslash{}(0.400 \textbackslash{}, \textbackslash{}text\{m\}\textbackslash{}) and the \textbackslash{}(y\textbackslash{})-axis with \textbackslash{}(0.300 \textbackslash{}, \textbackslash{}text\{m\}\textbackslash{}) to show distances. The arrows indicate directions and possibly movement or forces acting on the particles.

\#\# Dataset metadata

Magnetism; Physical Model Grounding Reasoning, Spatial Relation Reasoning

\paragraph{Recorded answer/context.}
B: B: (7.49 \textbackslash{}times 10\textasciicircum{}\{-8\}N)\textbackslash{}hat\{\textbackslash{}mathbf\{j\}\}

\paragraph{Figure/image path.}
/root/proposal\_for\_physic/science-mango/phyx\_mini\_run/phyx\_data/test\_image/963.png

\paragraph{Formalization target.}
create a compiling Lean file with sorry bodies at `PhyXMiniProblems/problem\_phyx\_mini\_0963.lean`. The Lean declarations must preserve the physical quantities, dimensions or dimensional roles, figure labels, governing-law hypotheses, and final relation expressed by this problem.
Use Mathlib/Physlib names found through LeanExplore where available. If a physics API is missing, introduce faithful local abstractions rather than scalar placeholder aliases.

\begin{theorem}[Physics formalization target]
\label{thm:physics:phyx_mini_0963:target}
\lean{PhyXMiniProblems.ProblemPhyXMini0963.problem_phyx_mini_0963}
\uses{def:physics:phyx-mini-0963:phyxminiproblems-problemphyxmini0963-forcevectorinnewtons, def:physics:phyx-mini-0963:phyxminiproblems-problemphyxmini0963-axisvector, def:physics:phyx-mini-0963:phyxminiproblems-problemphyxmini0963-pointsinaxisdirection, def:physics:phyx-mini-0963:phyxminiproblems-problemphyxmini0963-movingpointchargeforcesetup, def:physics:phyx-mini-0963:phyxminiproblems-problemphyxmini0963-matchesmovingpointchargescenario, def:physics:phyx-mini-0963:phyxminiproblems-problemphyxmini0963-matcheswrittenchargeandspeeddata, def:physics:phyx-mini-0963:phyxminiproblems-problemphyxmini0963-matchessuppliedmovingchargefigure, def:physics:phyx-mini-0963:phyxminiproblems-problemphyxmini0963-hasdepictedinstantaneouskinematics, def:physics:phyx-mini-0963:phyxminiproblems-problemphyxmini0963-hasphysicalmovingchargeparameters, def:physics:phyx-mini-0963:phyxminiproblems-problemphyxmini0963-usesstandardvacuumpermeability, def:physics:phyx-mini-0963:phyxminiproblems-problemphyxmini0963-satisfiesmovingpointchargemagneticfieldlaw, def:physics:phyx-mini-0963:phyxminiproblems-problemphyxmini0963-satisfiesmagneticlorentzforcelaw, def:physics:phyx-mini-0963:phyxminiproblems-problemphyxmini0963-displayedforcevectorinnewtons, def:physics:phyx-mini-0963:phyxminiproblems-problemphyxmini0963-roundstodisplayedforceprecision, def:physics:phyx-mini-0963:phyxminiproblems-problemphyxmini0963-isuniqueclosestdisplayedforceanswer}
The assigned autoformalize agent should translate this physics problem into Lean declarations in the covered file, with theorem and lemma proof bodies written as `by sorry`.
\end{theorem}
\begin{proof}
This is an autoformalization task, not a proof task. Produce statements that can later be proved by the physics prover without weakening the physical model.
\end{proof}

% --- Archon declaration topology begin ---
\section{Declaration topology}
\begin{definition}[force Dimension]
  \label{def:physics:phyx-mini-0963:phyxminiproblems-problemphyxmini0963-forcedimension}
  \lean{PhyXMiniProblems.ProblemPhyXMini0963.forceDimension}
  The physical dimension M L T\textasciicircum{}-2 of force
\end{definition}
\begin{proof}
  This is the declared model definition of force Dimension.
\end{proof}
\begin{definition}[Spatial Vector]
  \label{def:physics:phyx-mini-0963:phyxminiproblems-problemphyxmini0963-spatialvector}
  \lean{PhyXMiniProblems.ProblemPhyXMini0963.SpatialVector}
  Three-dimensional coherent-SI spatial vector readout
\end{definition}
\begin{proof}
  This is the declared model definition of Spatial Vector.
\end{proof}
\begin{definition}[Signed Charge Quantity]
  \label{def:physics:phyx-mini-0963:phyxminiproblems-problemphyxmini0963-signedchargequantity}
  \lean{PhyXMiniProblems.ProblemPhyXMini0963.SignedChargeQuantity}
  A signed, unit-independent electric charge
\end{definition}
\begin{proof}
  This is the declared model definition of Signed Charge Quantity.
\end{proof}
\begin{definition}[Speed Quantity]
  \label{def:physics:phyx-mini-0963:phyxminiproblems-problemphyxmini0963-speedquantity}
  \lean{PhyXMiniProblems.ProblemPhyXMini0963.SpeedQuantity}
  A nonnegative, unit-independent speed
\end{definition}
\begin{proof}
  This is the declared model definition of Speed Quantity.
\end{proof}
\begin{definition}[Force Vector Quantity]
  \label{def:physics:phyx-mini-0963:phyxminiproblems-problemphyxmini0963-forcevectorquantity}
  \lean{PhyXMiniProblems.ProblemPhyXMini0963.ForceVectorQuantity}
  \uses{def:physics:phyx-mini-0963:phyxminiproblems-problemphyxmini0963-forcedimension, def:physics:phyx-mini-0963:phyxminiproblems-problemphyxmini0963-spatialvector}
  A unit-independent three-dimensional force vector
\end{definition}
\begin{proof}
  This is the declared model definition of Force Vector Quantity.
\end{proof}
\begin{definition}[charge In Coulombs]
  \label{def:physics:phyx-mini-0963:phyxminiproblems-problemphyxmini0963-chargeincoulombs}
  \lean{PhyXMiniProblems.ProblemPhyXMini0963.chargeInCoulombs}
  \uses{def:physics:phyx-mini-0963:phyxminiproblems-problemphyxmini0963-signedchargequantity}
  Coherent-SI charge readout, in coulombs
\end{definition}
\begin{proof}
  This is the declared model definition of charge In Coulombs.
\end{proof}
\begin{definition}[charge In Microcoulombs]
  \label{def:physics:phyx-mini-0963:phyxminiproblems-problemphyxmini0963-chargeinmicrocoulombs}
  \lean{PhyXMiniProblems.ProblemPhyXMini0963.chargeInMicrocoulombs}
  \uses{def:physics:phyx-mini-0963:phyxminiproblems-problemphyxmini0963-signedchargequantity, def:physics:phyx-mini-0963:phyxminiproblems-problemphyxmini0963-chargeincoulombs}
  Microcoulomb readout used by the written charge data
\end{definition}
\begin{proof}
  This is the declared model definition of charge In Microcoulombs.
\end{proof}
\begin{definition}[speed In Meters Per Second]
  \label{def:physics:phyx-mini-0963:phyxminiproblems-problemphyxmini0963-speedinmeterspersecond}
  \lean{PhyXMiniProblems.ProblemPhyXMini0963.speedInMetersPerSecond}
  \uses{def:physics:phyx-mini-0963:phyxminiproblems-problemphyxmini0963-speedquantity}
  Coherent-SI speed readout, in metres per second
\end{definition}
\begin{proof}
  This is the declared model definition of speed In Meters Per Second.
\end{proof}
\begin{definition}[force Vector In Newtons]
  \label{def:physics:phyx-mini-0963:phyxminiproblems-problemphyxmini0963-forcevectorinnewtons}
  \lean{PhyXMiniProblems.ProblemPhyXMini0963.forceVectorInNewtons}
  \uses{def:physics:phyx-mini-0963:phyxminiproblems-problemphyxmini0963-spatialvector, def:physics:phyx-mini-0963:phyxminiproblems-problemphyxmini0963-forcevectorquantity}
  Coherent-SI force-vector readout, in newtons
\end{definition}
\begin{proof}
  This is the declared model definition of force Vector In Newtons.
\end{proof}
\begin{definition}[Particle Label]
  \label{def:physics:phyx-mini-0963:phyxminiproblems-problemphyxmini0963-particlelabel}
  \lean{PhyXMiniProblems.ProblemPhyXMini0963.ParticleLabel}
  The two particles, retaining the prime on the source charge as a name
\end{definition}
\begin{proof}
  This is the declared model definition of Particle Label.
\end{proof}
\begin{definition}[Coordinate Axis]
  \label{def:physics:phyx-mini-0963:phyxminiproblems-problemphyxmini0963-coordinateaxis}
  \lean{PhyXMiniProblems.ProblemPhyXMini0963.CoordinateAxis}
  The three Cartesian axes; the raster explicitly draws only x and y
\end{definition}
\begin{proof}
  This is the declared model definition of Coordinate Axis.
\end{proof}
\begin{definition}[Axis Direction]
  \label{def:physics:phyx-mini-0963:phyxminiproblems-problemphyxmini0963-axisdirection}
  \lean{PhyXMiniProblems.ProblemPhyXMini0963.AxisDirection}
  Signed directions along a Cartesian axis
\end{definition}
\begin{proof}
  This is the declared model definition of Axis Direction.
\end{proof}
\begin{definition}[Charge Sign Glyph]
  \label{def:physics:phyx-mini-0963:phyxminiproblems-problemphyxmini0963-chargesignglyph}
  \lean{PhyXMiniProblems.ProblemPhyXMini0963.ChargeSignGlyph}
  Plus/minus glyph shown inside a particle marker
\end{definition}
\begin{proof}
  This is the declared model definition of Charge Sign Glyph.
\end{proof}
\begin{definition}[Particle Marker Color]
  \label{def:physics:phyx-mini-0963:phyxminiproblems-problemphyxmini0963-particlemarkercolor}
  \lean{PhyXMiniProblems.ProblemPhyXMini0963.ParticleMarkerColor}
  Marker colours distinguished in the supplied raster
\end{definition}
\begin{proof}
  This is the declared model definition of Particle Marker Color.
\end{proof}
\begin{definition}[axis Vector]
  \label{def:physics:phyx-mini-0963:phyxminiproblems-problemphyxmini0963-axisvector}
  \lean{PhyXMiniProblems.ProblemPhyXMini0963.axisVector}
  \uses{def:physics:phyx-mini-0963:phyxminiproblems-problemphyxmini0963-spatialvector, def:physics:phyx-mini-0963:phyxminiproblems-problemphyxmini0963-coordinateaxis}
  Standard coherent-SI unit vector associated with an axis
\end{definition}
\begin{proof}
  This is the declared model definition of axis Vector.
\end{proof}
\begin{definition}[axis Direction Vector]
  \label{def:physics:phyx-mini-0963:phyxminiproblems-problemphyxmini0963-axisdirectionvector}
  \lean{PhyXMiniProblems.ProblemPhyXMini0963.axisDirectionVector}
  \uses{def:physics:phyx-mini-0963:phyxminiproblems-problemphyxmini0963-spatialvector, def:physics:phyx-mini-0963:phyxminiproblems-problemphyxmini0963-axisdirection, def:physics:phyx-mini-0963:phyxminiproblems-problemphyxmini0963-axisvector}
  Unit vector associated with an oriented Cartesian direction
\end{definition}
\begin{proof}
  This is the declared model definition of axis Direction Vector.
\end{proof}
\begin{definition}[spatial Cross]
  \label{def:physics:phyx-mini-0963:phyxminiproblems-problemphyxmini0963-spatialcross}
  \lean{PhyXMiniProblems.ProblemPhyXMini0963.spatialCross}
  \uses{def:physics:phyx-mini-0963:phyxminiproblems-problemphyxmini0963-spatialvector}
  Mathlib's right-handed cross product transported to Euclidean space
\end{definition}
\begin{proof}
  This is the declared model definition of spatial Cross.
\end{proof}
\begin{definition}[space Point Of Meters]
  \label{def:physics:phyx-mini-0963:phyxminiproblems-problemphyxmini0963-spacepointofmeters}
  \lean{PhyXMiniProblems.ProblemPhyXMini0963.spacePointOfMeters}
  \uses{def:physics:phyx-mini-0963:phyxminiproblems-problemphyxmini0963-spatialvector}
  Interpret a metre-coordinate vector as a point of Physlib space
\end{definition}
\begin{proof}
  This is the declared model definition of space Point Of Meters.
\end{proof}
\begin{definition}[position Vector In Meters]
  \label{def:physics:phyx-mini-0963:phyxminiproblems-problemphyxmini0963-positionvectorinmeters}
  \lean{PhyXMiniProblems.ProblemPhyXMini0963.positionVectorInMeters}
  \uses{def:physics:phyx-mini-0963:phyxminiproblems-problemphyxmini0963-spatialvector}
  Read the coordinates of a Physlib point in the chosen coherent-SI chart
\end{definition}
\begin{proof}
  This is the declared model definition of position Vector In Meters.
\end{proof}
\begin{definition}[Points In Axis Direction]
  \label{def:physics:phyx-mini-0963:phyxminiproblems-problemphyxmini0963-pointsinaxisdirection}
  \lean{PhyXMiniProblems.ProblemPhyXMini0963.PointsInAxisDirection}
  \uses{def:physics:phyx-mini-0963:phyxminiproblems-problemphyxmini0963-spatialvector, def:physics:phyx-mini-0963:phyxminiproblems-problemphyxmini0963-axisdirection, def:physics:phyx-mini-0963:phyxminiproblems-problemphyxmini0963-axisdirectionvector}
  A nonzero vector points along an oriented axis when it is a positive multiple of the corresponding unit vector
\end{definition}
\begin{proof}
  This is the declared model definition of Points In Axis Direction.
\end{proof}
\begin{definition}[Moving Point Charges Figure]
  \label{def:physics:phyx-mini-0963:phyxminiproblems-problemphyxmini0963-movingpointchargesfigure}
  \lean{PhyXMiniProblems.ProblemPhyXMini0963.MovingPointChargesFigure}
  \uses{def:physics:phyx-mini-0963:phyxminiproblems-problemphyxmini0963-particlelabel, def:physics:phyx-mini-0963:phyxminiproblems-problemphyxmini0963-coordinateaxis, def:physics:phyx-mini-0963:phyxminiproblems-problemphyxmini0963-axisdirection, def:physics:phyx-mini-0963:phyxminiproblems-problemphyxmini0963-chargesignglyph, def:physics:phyx-mini-0963:phyxminiproblems-problemphyxmini0963-particlemarkercolor}
  Literal marker and arrow data visible in primary image 963.png
\end{definition}
\begin{proof}
  This is the declared model definition of Moving Point Charges Figure.
\end{proof}
\begin{definition}[expected Particle Label]
  \label{def:physics:phyx-mini-0963:phyxminiproblems-problemphyxmini0963-expectedparticlelabel}
  \lean{PhyXMiniProblems.ProblemPhyXMini0963.expectedParticleLabel}
  \uses{def:physics:phyx-mini-0963:phyxminiproblems-problemphyxmini0963-particlelabel}
  Particle labels transcribed from the raster
\end{definition}
\begin{proof}
  This is the declared model definition of expected Particle Label.
\end{proof}
\begin{definition}[expected Velocity Label]
  \label{def:physics:phyx-mini-0963:phyxminiproblems-problemphyxmini0963-expectedvelocitylabel}
  \lean{PhyXMiniProblems.ProblemPhyXMini0963.expectedVelocityLabel}
  \uses{def:physics:phyx-mini-0963:phyxminiproblems-problemphyxmini0963-particlelabel}
  Velocity labels transcribed from the raster
\end{definition}
\begin{proof}
  This is the declared model definition of expected Velocity Label.
\end{proof}
\begin{definition}[expected Marker Color]
  \label{def:physics:phyx-mini-0963:phyxminiproblems-problemphyxmini0963-expectedmarkercolor}
  \lean{PhyXMiniProblems.ProblemPhyXMini0963.expectedMarkerColor}
  \uses{def:physics:phyx-mini-0963:phyxminiproblems-problemphyxmini0963-particlelabel, def:physics:phyx-mini-0963:phyxminiproblems-problemphyxmini0963-particlemarkercolor}
  Marker colour transcribed from the raster
\end{definition}
\begin{proof}
  This is the declared model definition of expected Marker Color.
\end{proof}
\begin{definition}[expected Sign Glyph]
  \label{def:physics:phyx-mini-0963:phyxminiproblems-problemphyxmini0963-expectedsignglyph}
  \lean{PhyXMiniProblems.ProblemPhyXMini0963.expectedSignGlyph}
  \uses{def:physics:phyx-mini-0963:phyxminiproblems-problemphyxmini0963-particlelabel, def:physics:phyx-mini-0963:phyxminiproblems-problemphyxmini0963-chargesignglyph}
  Charge-sign glyph transcribed from the raster
\end{definition}
\begin{proof}
  This is the declared model definition of expected Sign Glyph.
\end{proof}
\begin{definition}[Moving Point Charge Force Setup]
  \label{def:physics:phyx-mini-0963:phyxminiproblems-problemphyxmini0963-movingpointchargeforcesetup}
  \lean{PhyXMiniProblems.ProblemPhyXMini0963.MovingPointChargeForceSetup}
  \uses{def:physics:phyx-mini-0963:phyxminiproblems-problemphyxmini0963-spatialvector, def:physics:phyx-mini-0963:phyxminiproblems-problemphyxmini0963-signedchargequantity, def:physics:phyx-mini-0963:phyxminiproblems-problemphyxmini0963-speedquantity, def:physics:phyx-mini-0963:phyxminiproblems-problemphyxmini0963-forcevectorquantity, def:physics:phyx-mini-0963:phyxminiproblems-problemphyxmini0963-particlelabel, def:physics:phyx-mini-0963:phyxminiproblems-problemphyxmini0963-movingpointchargesfigure}
  The magnetic field and force are independent observables. In particular, neither is defined from the requested numerical answer or from an answer choice
\end{definition}
\begin{proof}
  This is the declared model definition of Moving Point Charge Force Setup.
\end{proof}
\begin{definition}[velocity Vector In Meters Per Second]
  \label{def:physics:phyx-mini-0963:phyxminiproblems-problemphyxmini0963-velocityvectorinmeterspersecond}
  \lean{PhyXMiniProblems.ProblemPhyXMini0963.velocityVectorInMetersPerSecond}
  \uses{def:physics:phyx-mini-0963:phyxminiproblems-problemphyxmini0963-spatialvector, def:physics:phyx-mini-0963:phyxminiproblems-problemphyxmini0963-speedinmeterspersecond, def:physics:phyx-mini-0963:phyxminiproblems-problemphyxmini0963-particlelabel, def:physics:phyx-mini-0963:phyxminiproblems-problemphyxmini0963-movingpointchargeforcesetup}
  Instantaneous velocity vector of a particle, in metres per second
\end{definition}
\begin{proof}
  This is the declared model definition of velocity Vector In Meters Per Second.
\end{proof}
\begin{definition}[displacement From QPrime To Position In Meters]
  \label{def:physics:phyx-mini-0963:phyxminiproblems-problemphyxmini0963-displacementfromqprimetopositioninmeters}
  \lean{PhyXMiniProblems.ProblemPhyXMini0963.displacementFromQPrimeToPositionInMeters}
  \uses{def:physics:phyx-mini-0963:phyxminiproblems-problemphyxmini0963-spatialvector, def:physics:phyx-mini-0963:phyxminiproblems-problemphyxmini0963-positionvectorinmeters, def:physics:phyx-mini-0963:phyxminiproblems-problemphyxmini0963-movingpointchargeforcesetup}
  Displacement from q' to an arbitrary observation point, in metres
\end{definition}
\begin{proof}
  This is the declared model definition of displacement From QPrime To Position In Meters.
\end{proof}
\begin{definition}[displacement From QPrime To QIn Meters]
  \label{def:physics:phyx-mini-0963:phyxminiproblems-problemphyxmini0963-displacementfromqprimetoqinmeters}
  \lean{PhyXMiniProblems.ProblemPhyXMini0963.displacementFromQPrimeToQInMeters}
  \uses{def:physics:phyx-mini-0963:phyxminiproblems-problemphyxmini0963-spatialvector, def:physics:phyx-mini-0963:phyxminiproblems-problemphyxmini0963-movingpointchargeforcesetup, def:physics:phyx-mini-0963:phyxminiproblems-problemphyxmini0963-displacementfromqprimetopositioninmeters}
  Displacement from q' to q at the pictured instant, in metres
\end{definition}
\begin{proof}
  This is the declared model definition of displacement From QPrime To QIn Meters.
\end{proof}
\begin{definition}[magnetic Field At QIn Teslas]
  \label{def:physics:phyx-mini-0963:phyxminiproblems-problemphyxmini0963-magneticfieldatqinteslas}
  \lean{PhyXMiniProblems.ProblemPhyXMini0963.magneticFieldAtQInTeslas}
  \uses{def:physics:phyx-mini-0963:phyxminiproblems-problemphyxmini0963-spatialvector, def:physics:phyx-mini-0963:phyxminiproblems-problemphyxmini0963-movingpointchargeforcesetup}
  Magnetic field of q' evaluated at q, read in teslas
\end{definition}
\begin{proof}
  This is the declared model definition of magnetic Field At QIn Teslas.
\end{proof}
\begin{definition}[Matches Moving Point Charge Scenario]
  \label{def:physics:phyx-mini-0963:phyxminiproblems-problemphyxmini0963-matchesmovingpointchargescenario}
  \lean{PhyXMiniProblems.ProblemPhyXMini0963.MatchesMovingPointChargeScenario}
  \uses{def:physics:phyx-mini-0963:phyxminiproblems-problemphyxmini0963-movingpointchargeforcesetup}
  The prose treats both objects as point charges
\end{definition}
\begin{proof}
  This is the declared model definition of Matches Moving Point Charge Scenario.
\end{proof}
\begin{definition}[Matches Written Charge And Speed Data]
  \label{def:physics:phyx-mini-0963:phyxminiproblems-problemphyxmini0963-matcheswrittenchargeandspeeddata}
  \lean{PhyXMiniProblems.ProblemPhyXMini0963.MatchesWrittenChargeAndSpeedData}
  \uses{def:physics:phyx-mini-0963:phyxminiproblems-problemphyxmini0963-chargeinmicrocoulombs, def:physics:phyx-mini-0963:phyxminiproblems-problemphyxmini0963-speedinmeterspersecond, def:physics:phyx-mini-0963:phyxminiproblems-problemphyxmini0963-movingpointchargeforcesetup}
  Signed charge and speed data stated in the problem
\end{definition}
\begin{proof}
  This is the declared model definition of Matches Written Charge And Speed Data.
\end{proof}
\begin{definition}[Matches Supplied Moving Charge Figure]
  \label{def:physics:phyx-mini-0963:phyxminiproblems-problemphyxmini0963-matchessuppliedmovingchargefigure}
  \lean{PhyXMiniProblems.ProblemPhyXMini0963.MatchesSuppliedMovingChargeFigure}
  \uses{def:physics:phyx-mini-0963:phyxminiproblems-problemphyxmini0963-expectedparticlelabel, def:physics:phyx-mini-0963:phyxminiproblems-problemphyxmini0963-expectedvelocitylabel, def:physics:phyx-mini-0963:phyxminiproblems-problemphyxmini0963-expectedmarkercolor, def:physics:phyx-mini-0963:phyxminiproblems-problemphyxmini0963-expectedsignglyph, def:physics:phyx-mini-0963:phyxminiproblems-problemphyxmini0963-movingpointchargeforcesetup}
  Literal axes, labels, colours, positions, signs, and arrow directions from the primary raster. This predicate contains no force information
\end{definition}
\begin{proof}
  This is the declared model definition of Matches Supplied Moving Charge Figure.
\end{proof}
\begin{definition}[Has Depicted Instantaneous Kinematics]
  \label{def:physics:phyx-mini-0963:phyxminiproblems-problemphyxmini0963-hasdepictedinstantaneouskinematics}
  \lean{PhyXMiniProblems.ProblemPhyXMini0963.HasDepictedInstantaneousKinematics}
  \uses{def:physics:phyx-mini-0963:phyxminiproblems-problemphyxmini0963-axisvector, def:physics:phyx-mini-0963:phyxminiproblems-problemphyxmini0963-axisdirectionvector, def:physics:phyx-mini-0963:phyxminiproblems-problemphyxmini0963-spacepointofmeters, def:physics:phyx-mini-0963:phyxminiproblems-problemphyxmini0963-movingpointchargeforcesetup}
  Metric interpretation of the two position labels and the two arrow directions. The unpictured z coordinate is zero
\end{definition}
\begin{proof}
  This is the declared model definition of Has Depicted Instantaneous Kinematics.
\end{proof}
\begin{definition}[Has Physical Moving Charge Parameters]
  \label{def:physics:phyx-mini-0963:phyxminiproblems-problemphyxmini0963-hasphysicalmovingchargeparameters}
  \lean{PhyXMiniProblems.ProblemPhyXMini0963.HasPhysicalMovingChargeParameters}
  \uses{def:physics:phyx-mini-0963:phyxminiproblems-problemphyxmini0963-speedinmeterspersecond, def:physics:phyx-mini-0963:phyxminiproblems-problemphyxmini0963-movingpointchargeforcesetup, def:physics:phyx-mini-0963:phyxminiproblems-problemphyxmini0963-displacementfromqprimetoqinmeters}
  Positivity and nondegeneracy implicit in two distinct moving particles
\end{definition}
\begin{proof}
  This is the declared model definition of Has Physical Moving Charge Parameters.
\end{proof}
\begin{definition}[Uses Standard Vacuum Permeability]
  \label{def:physics:phyx-mini-0963:phyxminiproblems-problemphyxmini0963-usesstandardvacuumpermeability}
  \lean{PhyXMiniProblems.ProblemPhyXMini0963.UsesStandardVacuumPermeability}
  \uses{def:physics:phyx-mini-0963:phyxminiproblems-problemphyxmini0963-movingpointchargeforcesetup}
  Standard coherent-SI value μ₀ = 4π * 10\textasciicircum{}-7 used by the model
\end{definition}
\begin{proof}
  This is the declared model definition of Uses Standard Vacuum Permeability.
\end{proof}
\begin{definition}[Satisfies Moving Point Charge Magnetic Field Law]
  \label{def:physics:phyx-mini-0963:phyxminiproblems-problemphyxmini0963-satisfiesmovingpointchargemagneticfieldlaw}
  \lean{PhyXMiniProblems.ProblemPhyXMini0963.SatisfiesMovingPointChargeMagneticFieldLaw}
  \uses{def:physics:phyx-mini-0963:phyxminiproblems-problemphyxmini0963-chargeincoulombs, def:physics:phyx-mini-0963:phyxminiproblems-problemphyxmini0963-spatialcross, def:physics:phyx-mini-0963:phyxminiproblems-problemphyxmini0963-movingpointchargeforcesetup, def:physics:phyx-mini-0963:phyxminiproblems-problemphyxmini0963-velocityvectorinmeterspersecond, def:physics:phyx-mini-0963:phyxminiproblems-problemphyxmini0963-displacementfromqprimetopositioninmeters}
  Quasistatic magnetic field of the moving source point charge q': B(x) = μ₀ q' / (4π |r|\textasciicircum{}3) (v' × r), where r = x - x\_\{q'\}. This is a general field law away from the source. It does not state the requested force or its answer-choice value
\end{definition}
\begin{proof}
  This is the declared model definition of Satisfies Moving Point Charge Magnetic Field Law.
\end{proof}
\begin{definition}[Satisfies Magnetic Lorentz Force Law]
  \label{def:physics:phyx-mini-0963:phyxminiproblems-problemphyxmini0963-satisfiesmagneticlorentzforcelaw}
  \lean{PhyXMiniProblems.ProblemPhyXMini0963.SatisfiesMagneticLorentzForceLaw}
  \uses{def:physics:phyx-mini-0963:phyxminiproblems-problemphyxmini0963-chargeincoulombs, def:physics:phyx-mini-0963:phyxminiproblems-problemphyxmini0963-forcevectorinnewtons, def:physics:phyx-mini-0963:phyxminiproblems-problemphyxmini0963-spatialcross, def:physics:phyx-mini-0963:phyxminiproblems-problemphyxmini0963-movingpointchargeforcesetup, def:physics:phyx-mini-0963:phyxminiproblems-problemphyxmini0963-velocityvectorinmeterspersecond, def:physics:phyx-mini-0963:phyxminiproblems-problemphyxmini0963-magneticfieldatqinteslas}
  Magnetic part of the Lorentz-force law on q, F = q (v × B). The force observable on the left is not otherwise constrained to the requested value
\end{definition}
\begin{proof}
  This is the declared model definition of Satisfies Magnetic Lorentz Force Law.
\end{proof}
\begin{lemma}[displacement From QPrime To Q exact]
  \label{lem:physics:phyx-mini-0963:phyxminiproblems-problemphyxmini0963-displacementfromqprimetoq-exact}
  \lean{PhyXMiniProblems.ProblemPhyXMini0963.displacementFromQPrimeToQ_exact}
  \uses{def:physics:phyx-mini-0963:phyxminiproblems-problemphyxmini0963-axisvector, def:physics:phyx-mini-0963:phyxminiproblems-problemphyxmini0963-movingpointchargeforcesetup, def:physics:phyx-mini-0963:phyxminiproblems-problemphyxmini0963-displacementfromqprimetoqinmeters, def:physics:phyx-mini-0963:phyxminiproblems-problemphyxmini0963-hasdepictedinstantaneouskinematics}
  The depicted geometry gives the familiar 3-4-5 displacement vector from q' to q
\end{lemma}
\begin{proof}
  The conclusion follows from the governing relations and figure readouts named in the statement of displacement From QPrime To Q exact.
\end{proof}
\begin{lemma}[magnetic Field At Q exact]
  \label{lem:physics:phyx-mini-0963:phyxminiproblems-problemphyxmini0963-magneticfieldatq-exact}
  \lean{PhyXMiniProblems.ProblemPhyXMini0963.magneticFieldAtQ_exact}
  \uses{def:physics:phyx-mini-0963:phyxminiproblems-problemphyxmini0963-axisvector, def:physics:phyx-mini-0963:phyxminiproblems-problemphyxmini0963-movingpointchargeforcesetup, def:physics:phyx-mini-0963:phyxminiproblems-problemphyxmini0963-magneticfieldatqinteslas, def:physics:phyx-mini-0963:phyxminiproblems-problemphyxmini0963-matcheswrittenchargeandspeeddata, def:physics:phyx-mini-0963:phyxminiproblems-problemphyxmini0963-hasdepictedinstantaneouskinematics, def:physics:phyx-mini-0963:phyxminiproblems-problemphyxmini0963-hasphysicalmovingchargeparameters, def:physics:phyx-mini-0963:phyxminiproblems-problemphyxmini0963-usesstandardvacuumpermeability, def:physics:phyx-mini-0963:phyxminiproblems-problemphyxmini0963-satisfiesmovingpointchargemagneticfieldlaw}
  The source charge produces a -z magnetic field of 1.04 * 10\textasciicircum{}-7 T at q
\end{lemma}
\begin{proof}
  The conclusion follows from the governing relations and figure readouts named in the statement of magnetic Field At Q exact.
\end{proof}
\begin{definition}[Answer Choice]
  \label{def:physics:phyx-mini-0963:phyxminiproblems-problemphyxmini0963-answerchoice}
  \lean{PhyXMiniProblems.ProblemPhyXMini0963.AnswerChoice}
  Labels attached to the four displayed multiple-choice answers
\end{definition}
\begin{proof}
  This is the declared model definition of Answer Choice.
\end{proof}
\begin{definition}[displayed Force Vector In Newtons]
  \label{def:physics:phyx-mini-0963:phyxminiproblems-problemphyxmini0963-displayedforcevectorinnewtons}
  \lean{PhyXMiniProblems.ProblemPhyXMini0963.displayedForceVectorInNewtons}
  \uses{def:physics:phyx-mini-0963:phyxminiproblems-problemphyxmini0963-spatialvector, def:physics:phyx-mini-0963:phyxminiproblems-problemphyxmini0963-axisvector, def:physics:phyx-mini-0963:phyxminiproblems-problemphyxmini0963-answerchoice}
  Force vector in newtons printed beside each answer label
\end{definition}
\begin{proof}
  This is the declared model definition of displayed Force Vector In Newtons.
\end{proof}
\begin{definition}[recorded Dataset Answer]
  \label{def:physics:phyx-mini-0963:phyxminiproblems-problemphyxmini0963-recordeddatasetanswer}
  \lean{PhyXMiniProblems.ProblemPhyXMini0963.recordedDatasetAnswer}
  \uses{def:physics:phyx-mini-0963:phyxminiproblems-problemphyxmini0963-answerchoice}
  The answer label recorded in the source dataset
\end{definition}
\begin{proof}
  This is the declared model definition of recorded Dataset Answer.
\end{proof}
\begin{definition}[Rounds To Displayed Force Precision]
  \label{def:physics:phyx-mini-0963:phyxminiproblems-problemphyxmini0963-roundstodisplayedforceprecision}
  \lean{PhyXMiniProblems.ProblemPhyXMini0963.RoundsToDisplayedForcePrecision}
  \uses{def:physics:phyx-mini-0963:phyxminiproblems-problemphyxmini0963-spatialvector}
  The vector agrees with a displayed vector to the 0.01 * 10\textasciicircum{}-8 N precision used in choice B
\end{definition}
\begin{proof}
  This is the declared model definition of Rounds To Displayed Force Precision.
\end{proof}
\begin{definition}[Is Unique Closest Displayed Force Answer]
  \label{def:physics:phyx-mini-0963:phyxminiproblems-problemphyxmini0963-isuniqueclosestdisplayedforceanswer}
  \lean{PhyXMiniProblems.ProblemPhyXMini0963.IsUniqueClosestDisplayedForceAnswer}
  \uses{def:physics:phyx-mini-0963:phyxminiproblems-problemphyxmini0963-spatialvector, def:physics:phyx-mini-0963:phyxminiproblems-problemphyxmini0963-answerchoice, def:physics:phyx-mini-0963:phyxminiproblems-problemphyxmini0963-displayedforcevectorinnewtons}
  A displayed answer is uniquely closest to the actual force vector
\end{definition}
\begin{proof}
  This is the declared model definition of Is Unique Closest Displayed Force Answer.
\end{proof}
% --- Archon declaration topology end ---
% --- Archon physics formalization source end ---
